% !PRINTED 564
\weberpage{564}
Wir wollen hier nicht den genauen Werth für $K_n$ einsetzen, sondern
einen einfacher auszudrückenden kleineren Werth, wodurch die
Ungleichung (11) um so mehr richtig wird.

Nach der Definition ist

\weq{17}{K_n = \int\!\int\!\cdots\!\int d\xi_1\, d\xi_2 \ldots d\xi_n,}

mit der Begrenzung

\weq{18}{\xi_1^2 + \xi_2^2 + \cdots + \xi_n^2 \leqq 1.}

Die Grenzgleichung (18) ist sicher erfüllt, wenn wir die $\xi_i$
auf die Intervalle

\weq{19}{-\frac{1}{\sqrt{n}} < \xi_i < \frac{1}{\sqrt{n}}}

beschränken, was geometrisch die Bedeutung hat, dass wir die Kugel
durch den eingeschriebenen Würfel ersetzen. Demnach ist

\[
K_n > \int\limits_{-\frac{1}{\sqrt n}}^{\frac{1}{\sqrt n}}
\int\limits_{-\frac{1}{\sqrt n}}^{\frac{1}{\sqrt n}}
\cdots
\int\limits_{-\frac{1}{\sqrt n}}^{\frac{1}{\sqrt n}}
d\xi_1\, d\xi_2 \cdots d\xi_n,
\]

und wenn diese Integration ausgeführt wird,

\weq{20}{K_n > \left(\frac{2}{\sqrt{n}}\right)^{n}.}

Die hierdurch bestimmte obere Grenze wird von $K_n$ niemals
erreicht, wenn wir von dem Falle $n = 1$ absehen.

Die Vergleichung von (20) mit (16) ergiebt einen sehr einfachen
Ausdruck für den oberen Grenzwerth von $M$, nämlich

\weq{21}{M < n \sqrt[n]{D},}

und $M$ ist (von dem Falle $n = 1$ abgesehen) wirklich kleiner, und
nicht gleich dem Werthe $n \sqrt[n]{D}$.

Dies zeigt sich klar, wenn man eine noch etwas engere Grenze für
$M$ aufsucht.

Die Bedingung (18) ist sicher dann erfüllt, wenn wir

\weq{22}{-\alpha < \xi_n < \alpha,\qquad \xi_1^2 + \xi_2^2 + \cdots + \xi_{n-1}^2 < 1 - \alpha^2}

setzen, wo $\alpha$ ein beliebiger positiver echter Bruch sein kann.
(Die geometrische Bedeutung hiervon ist die, dass man die Kugel
durch einen eingeschriebenen Cylinder von der Höhe $2\alpha$
ersetzt). Es ist also

\[
\begin{gathered}
K_n > 2\alpha \int\!\int\!\cdots\!\int d\xi_1\, d\xi_2 \ldots d\xi_{n-1}\\
\xi_1^2 + \xi_2^2 + \cdots + \xi_{n-1}^2 < 1 - \alpha^2.
\end{gathered}
\]
